\section{Interrogare l'Ontologia con SPARQL}
\label{sec:interrogazione}
\sloppy

Il valore degli assiomi si vede al momento dell'interrogazione: dopo il reasoning,
le conoscenze implicite sono triple esplicite del grafo e le query SPARQL possono
usarle direttamente, senza ricostruire la logica dello schema. Gli esempi seguenti
usano il prefisso \texttt{vg:} e mostrano quattro tipi di conoscenza derivata.

\subsection{Classi Inferite al Posto dei Filtri}

\begin{lstlisting}[style=sparql, caption={Giochi multiplayer per PlayStation con punteggio alto}, label={lst:q_classi}]
PREFIX vg: <http://www.videogame-ontology.org/ontology#>
SELECT ?name ?score WHERE {
  ?g a vg:HighlyRatedGame , vg:MultiplayerGame , vg:PlayStationGame ;
     vg:gameName ?name ;
     vg:metacriticScore ?score .
}
ORDER BY DESC(?score) LIMIT 10
\end{lstlisting}

Senza reasoning la stessa domanda richiederebbe un filtro sul punteggio
($\ge 85$), un join con le modalit\`a di gioco e un join con le piattaforme per
verificarne il tipo (ad esempio \texttt{PS5} \texttt{a} \texttt{PlayStationPlatform}). Con le classi definite (Sezione~\ref{sec:classi_definite})
la query si riduce a tre asserzioni di tipo.

\subsection{Relazioni Derivate da Property Chain}

\begin{lstlisting}[style=sparql, caption={Giochi dello stesso sviluppatore di Elden Ring}, label={lst:q_chain}]
PREFIX vg: <http://www.videogame-ontology.org/ontology#>
SELECT DISTINCT ?other WHERE {
  ?er vg:gameName "Elden Ring" ;
      vg:sharesDeveloperWith ?g .
  ?g  vg:gameName ?other .
  FILTER(?g != ?er)
}
\end{lstlisting}

\texttt{sharesDeveloperWith} non \`e mai asserita nei dati: deriva dalla catena
\texttt{developedBy} $\circ$ \texttt{developerOf} (Sezione~\ref{sec:chains}).
Il \texttt{FILTER} esclude la coppia riflessiva che la catena produce per ogni gioco.

\subsection{Transitivit\`a: Materializzazione o Property Path}

\begin{lstlisting}[style=sparql, caption={Tutti i predecessori di un gioco}, label={lst:q_trans}]
PREFIX vg: <http://www.videogame-ontology.org/ontology#>
SELECT ?pred WHERE {
  ?g vg:gameName "Dark Souls III" ;
     vg:sequelOf ?pred .      # con reasoning (chiusura materializzata)
  # senza reasoning:  ?g vg:sequelOf+ ?pred .
}
\end{lstlisting}

Poich\'e \texttt{sequelOf} \`e transitiva, dopo il reasoning il grafo contiene gi\`a
tutte le coppie (gioco, predecessore) a qualsiasi distanza. Lo stesso risultato si
pu\`o ottenere senza reasoning con il \emph{property path} SPARQL
\texttt{vg:sequelOf+}, che calcola la chiusura al momento della query. La differenza
\`e concettuale: nel primo caso la transitivit\`a \`e una propriet\`a del dominio
dichiarata nell'ontologia e vale per ogni applicazione che la usa; nel secondo \`e
una scelta di chi scrive la query.

\subsection{Navigazione Tramite Inverse}

\begin{lstlisting}[style=sparql, caption={Studi con pi\`u giochi premiati}, label={lst:q_inv}]
PREFIX vg: <http://www.videogame-ontology.org/ontology#>
SELECT ?studio (COUNT(DISTINCT ?g) AS ?n) WHERE {
  ?dev vg:developerOf ?g ;
       vg:developerName ?studio .
  ?g a vg:AwardWinningGame .
}
GROUP BY ?studio ORDER BY DESC(?n) LIMIT 10
\end{lstlisting}

I dati asseriscono solo \texttt{developedBy} (dal gioco allo studio): la query parte
dallo studio grazie all'inversa \texttt{developerOf} e usa la classe definita
\texttt{AwardWinningGame} invece di un join esplicito con i premi.

\subsection{Open World e Negazione}

Una domanda come \textless{}giochi senza sequel\textgreater{} mette in luce la differenza
tra l'ontologia e le query. In OWL vale la \textit{Open World Assumption}: l'assenza di
una tripla \texttt{sequelOf} non implica che il gioco non abbia predecessori, quindi un
reasoner DL non classificherebbe nessun gioco come \texttt{StandaloneGame} senza
ulteriori assiomi di chiusura. Il secondo livello di reasoning adotta invece
esplicitamente una \textit{Closed World Assumption} locale (\texttt{FILTER NOT EXISTS},
Sezione~\ref{sec:reasoning}): \texttt{StandaloneGame} va quindi letta come
\textless{}gioco per cui non \`e noto alcun predecessore\textgreater{}.

\fussy
